\section{Résultats relatifs à l'amendement du sol}
Comme le décrivent nos chapitres précédents, pour chaque amendement utilisé, nous avons des attentes spécifiques. Cette section décrit d’abord les résultats attendus selon la littérature scientifique avant de présenter les résultats réellement obtenus.

\subsection{Mécanisme d’action attendu des composants de l’amendement}
Les différents composants de l’amendement ont été associés en raison de leurs effets complémentaires attendus sur le sol et les cultures. \textbf{Le biochar} contribue à la rétention des nutriments et à l’amélioration des propriétés du sol, tandis que le \textbf{compost} favorise sa fertilité et sa capacité de rétention en eau. La \textbf{PCA} constitue une source de matière organique à décomposition lente et contribue aussi à l'aération du sol. Enfin, \textbf{l'EGB}, grâce aux micro-organismes bénéfiques qu’il contient, est susceptible de favoriser la nutrition des plantes et leur tolérance au stress salin.

Ces effets constituent les mécanismes attendus dont la manifestation doit être appréciée à partir des observations réalisées au cours de l’expérimentation.

\subsection{Résultats des tests au début de l'expérimentation}
La présente section est consacrée à l’interprétation des données obtenues selon trois axes complémentaires : 
l’évaluation du comportement du capteur, la caractérisation des composantes de l’amendement utilisé et l’analyse comparative de l’état du sol entre les plots T et les plots TS.

\subsubsection{Validation du comportement du capteur}
Avant d’utiliser les données du sol collectées sur le terrain, la cohérence de la réponse du capteur a été vérifiée à différents niveaux d’humidité. Un même échantillon de sol non amendé a été réparti dans trois récipients correspondant à trois niveaux d'humidité différents : \textbf{sol sec, moitié humide et sol mouillé}. Les mesures obtenues de ces pots ont été relevées et leurs valeurs sont présentées dans le tableau ci-dessous :

\begin{table}[H]
    \centering
    \caption{Résultats sur les échantillons des pots}
    \label{tab:collecte_pots}
    \begin{tabular}{|p{2.4cm}|p{1.6cm}|p{1.9cm}|p{0.8cm}|p{1.9cm}|p{1.7cm}|p{1.7cm}|p{1.7cm}|}
        \hline
        \textbf{État du sol} & \textbf{Humidité (\%)} & \textbf{Température (°C)} & \textbf{pH} & \textbf{EC (\textmu S/cm)} & \textbf{N (mg/kg)} & \textbf{P (mg/kg)} & \textbf{K (mg/kg)} \\
        \hline
        Sol sec & 0--5 & 32,7--32,9 & 7 & 0 & 0 & 0 & 0 \\
        \hline
        Sol moitié humide & 7,2--7,8 & 31,9 & 7 & 126--193 & 6--9 & 8--13 & 20--30 \\
        \hline
        Sol mouillé & 100 & 31,3 & 10 & 5451--5570 & 272--278 & 381--389 & 872--891 \\
        \hline
    \end{tabular}
\end{table}

Les résultats montrent une augmentation régulière de $EC$, du $N$, du $P$ et du $K$ avec l'humidité du sol. Cette tendance s’explique par le fait que l’eau favorise la dissolution et le déplacement des ions, facilitant ainsi leur détection par le capteur. À l’inverse, un sol sec limite fortement leur déplacement, ce qui explique les valeurs nulles observées dans le pot de sol sec. La valeur de 100\% d'humidité mesurée dans le sol mouillé s’explique par sa saturation uniforme contrairement aux plots au champ ou le drainage naturel limite une telle saturation. Ce test en pot constitue donc principalement une vérification de la cohérence de la réponse du capteur à des niveaux d'humidité différents.

\begin{figure}[H]
    \centering
    \imgstd{images/27_echantillons_de_sol.jpg}
    \caption{Échantillons de sol à différents niveaux d'humidité. }
    \label{fig:echantillons3}
\end{figure}

\subsubsection{Caractérisation des composantes de l’amendement}
L'amendement anti-sel est constitué de compost, de biochar, de PCA et de EGB. Afin d’identifier la contribution de chacun, un test d’hydratation individuel a été effectué avant leur incorporation au sol. Les résultats obtenus sont présentés dans le tableau ci-dessous :

\begin{table}[H]
    \centering
    \small
    \caption{Résultats sur les composants de l'amendement}
    \label{tab:releves_amendement}
    \begin{tabular}{|p{2.4cm}|p{1.6cm}|p{1.9cm}|p{0.9cm}|p{1.9cm}|p{1.5cm}|p{1.6cm}|p{1.6cm}|}
        \hline
        \textbf{Composants} & \textbf{Humidité (\%)} & \textbf{Température (°C)} & \textbf{pH} & \textbf{EC (\textmu S/cm)} & \textbf{N (mg/kg)} & \textbf{P (mg/kg)} & \textbf{K (mg/kg)} \\
        \hline
        Biofertilisant & 42,9--43,6 & 31 & 6,9 & 1394--1397 & 69 & 97 & 223 \\
        \hline
        Biochar & 75,3--75,6 & 30,7--30,8 & 6,8--6,9 & 1605--1659 & 80--82 & 112--116 & 256--265 \\
        \hline
        Compost & 100 & 30,4 & 6,9--7,6 & 6112--6208 & 305--310 & 427--434 & 977--993 \\
        \hline
        Coque d'arachide & 52,2--52,3 & 30--31,5 & 6,9 & 1080--1085 & 54 & 75 & 172--173 \\
        \hline
    \end{tabular}
\end{table}

Ces résultats révèlent que le compost présente les valeurs les plus élevées de $EC$, de $N$, de $P$ et de $K$, en raison de sa richesse en MO minéralisée. Malgré cette EC forte, son rôle dans la lutte contre la salinité repose principalement sur l'amélioration de la structure, de la rétention et du drainage du sol favorisant ainsi le lessivage des sels excédentaires. Le EGB et la PCA caractérisés par les $EC$ les plus faibles participent quant à eux à la régulation de la salinité dans le mélange. La figure ci-dessous montre les échantillons analysés.

\begin{figure}[H]
    \centering
    \imgstd{images/28_prelevement_intrant.jpg}
    \caption{Prélèvements des 4 éléments de l’amendement. }
    \label{fig:prelevement3}
\end{figure}

\subsubsection{Analyse comparative de l’état du sol : plots témoins et plots test}
Après mélange et incorporation des quatres composants dans les plots TS suivi d’un arrosage identique, un relevé a été réalisé la veille du semis (le 24 juin 2026) afin de comparer l'état initial du sol entre plot TS amendé et plot T non amendé.

\begin{table}[H]
    \centering
    \small
    \caption{Résultats de l'analyse comparative entre les plots test et témoins}
    \label{tab:analyse_comparative}
    \begin{tabular}{|p{1.9cm}|p{1.6cm}|p{1.9cm}|p{0.9cm}|p{1.9cm}|p{1.5cm}|p{1.6cm}|p{1.6cm}|}
        \hline
        \textbf{Plot} & \textbf{Humidité (\%)} & \textbf{Température (°C)} & \textbf{pH} & \textbf{EC (\textmu S/cm)} & \textbf{N (mg/kg)} & \textbf{P (mg/kg)} & \textbf{K (mg/kg)} \\
        \hline
        GOM T & 79,3 & 32,3 & 10 & 2282,2 & 113,7 & 159,3 & 364,7 \\
        \hline
        GOM TS & 100 & 30,3 & 8,2 & 4615,6 & 230,3 & 322,6 & 738 \\
        \hline
        NAV T & 71,5 & 31,7 & 10 & 3010,8 & 150 & 210,4 & 481,4 \\
        \hline
        NAV TS & 100 & 30,9 & 7,2 & 4681,9 & 233,5 & 327,3 & 748,6 \\
        \hline
        LAI T & 55,7 & 31 & 10 & 1374,8 & 68,2 & 95,7 & 219,5 \\
        \hline
        LAI TS & 100 & 31,4 & 7,3 & 4873,6 & 243,1 & 340,6 & 779,2 \\
        \hline
    \end{tabular}
\end{table}

L’analyse comparative met en évidence des différences nettes entre les plots T et les plots TS. Les plots TS présentent des valeurs de $EC$, de $N$, de $P$ et de $K$ environ deux à trois fois plus élevées traduisant un enrichissement du sol en liaison avec l’amendement. Le pH égal à 10 dans les plots T est ramené à 7,2-8.2 dans les plots TS, tandis que leur humidité atteint 100\% contre 55,7-79,3\% dans les plots T. Ces résultats suggèrent ainsi une amélioration initiale de la fertilité, de la rétention en eau du sol amendé et de la régulation de l'alcalinité.

\subsection{Évolution de l'état du sol entre plots tests et plots témoins (observations)} 
Quelques jours après l’incorporation des composants de l’amendement et les premiers arrosages, des différences marquées ont été observées entre les plots témoins et les plots tests pour nos trois cultures.

Sur les plots TS, le sol est resté meuble et friable, avec une infiltration satisfaisante tout au long du cycle. Sur les plots T, en particulier NAV T et LAI T, des fissures superficielles sont apparues et le sol semblait s’affaisser sur lui-même (Figure ci-dessous). 
L’eau d’irrigation ruisselait alors vers le côté en pente du plot au lieu de s’infiltrer uniformément.
\begin{figure}[H]
    \centering
    \imgstd{images/a_fissure_chap5.jpg}
    \caption{Illustration avec une image de plot (NAV T) qui s’affaisse. }
    \label{fig:fissure5}
\end{figure}

L’eau s’infiltrait rapidement sur les plots TS, alors qu’elle nécessitait au minimum cinq minutes sur les plots T (Figures ci-dessous). 
Après des pluies nocturnes intenses, l’eau était déjà fortement infiltrée dans les plots TS, tandis que les plots T restaient dans un état plus détérioré. 
Cette différence était plus visible sur NAV T et LAI T, semés en ligne, que sur GOM T, semé en poquets. Le mode de semis et d’autres facteurs propres à chaque plot ont pu contribuer à cette différence.
\begin{figure}[H]
    \centering
    \imgstd{images/b_gombo_chap5.jpg}
    \caption{Infiltration sur les plots GOM T et GOM TS 5 minutes après irrigation observée au douzième jour. }
    \label{fig:gombo5}
\end{figure}

\begin{figure}[H]
    \centering
    \imgstd{images/c_navet_chap5.jpg}
    \caption{Infiltration sur les plots NAV T et NAV TS 5 minutes après irrigation observée au douzième jour. }
    \label{fig:navet5}
\end{figure}

Pendant la journée, le sol des plots TS restait meuble, tandis que celui des plots T \textbf{durcissait} progressivement jusqu’à devenir aussi dur que du carrelage : il supportait le poids d’une personne sans \textbf{s’effondrer et était difficile à briser à la main} (Figure suivante). 
Les plots T ont nécessité un binage plus fréquent pour faciliter l’infiltration de l’eau et l’aération du sol.
\begin{figure}[H]
    \centering
    \imgstd{images/d_sol_dur_meuble_chap5.jpg}
    \caption{Illustration avec des images de sol T dur et de sol TS meuble. }
    \label{fig:dur_meuble5}
\end{figure}

Deux autres phénomènes ont été souvent observés sur les plots T et jamais sur les plots TS (Figure ci-dessous). Après arrosage, de légères taches blanches apparaissaient à la surface du sol et devenaient plus visibles après séchage, 
sous forme d’une efflorescence blanchâtre. Par ailleurs, quelques minutes après l’irrigation, l’eau, incolore au départ, devenait noirâtre à brunâtre. Enfin, le pH mesuré sur les plots témoins avant amendement était très proche de 10.
\begin{figure}[H]
    \centering
    \imgstd{images/e_signes_chap5.jpg}
    \caption{Illustration des signes observés uniquement sur les témoins : taches blanches indiquant la présence du sel 1, toxicité 2 et tache blanche + toxicité 3. }
    \label{fig:toxique5}
\end{figure}

\subsection{Interprétation et discussion sur les observations}
Les observations précédentes sont interprétées avec prudence, car aucune analyse du sol n’a permis de les confirmer directement.

\textbf{Interprétation des symptômes} : L’ensemble des symptômes observés sur les plots T (\textbf{fissuration, croûtage de surface, dureté à l’état sec, ruissellement préférentiel, taches blanches}) est compatible avec une situation de sol sodique : 
un taux élevé de sodium échangeable provoque la dispersion des argiles, la destruction de la structure du sol, la formation d’une croûte de battance qui limite l’infiltration de l’eau, et une fissuration lors du séchage \cite{niaz_organic_2022}. 
Cette \textbf{interprétation est cohérente} avec le pH très proche de 10 mesuré avant amendement. Les taches blanches pourraient indiquer une \textbf{remontée de sel par capillarité}, sans que leur nature ait été vérifiée.

\textbf{Effet attendu de l’amendement} : L’absence d’efflorescence sur les plots TS est compatible avec un effet du biochar sur la concentration en sels en surface, par fixation du sodium et entraînement en profondeur par lessivage. 
Le sol plus meuble et perméable sur les plots TS, contre un sol durci, fissuré et sujet à l’érosion sur les plots T, est cohérent avec \textbf{l’effet structurant du compost et de la PCA}. Le binage moins fréquent sur les plots TS suggère que cet effet s’est maintenu pendant le cycle. 
L’EGB agit surtout au niveau des racines, en favorisant l’absorption du potassium au détriment du sodium. Son impact n’étant pas directement observable sur le sol, son efficacité sera jugée à travers les données de croissance et de rendement.

\textbf{Observation non expliquée} : La coloration noirâtre de l’eau après irrigation, observée uniquement sur les plots T, reste inexpliquée. 
Elle pourrait être liée à la réduction de certains éléments (fer, manganèse) ou à un excès de sodium, mais cette hypothèse n’a pas été validée analytiquement et constitue une piste pour des \textbf{travaux futurs}.

\textbf{Limites} : Les composants de l’amendement ayant été appliqués ensemble, l’effet de chacun ne peut pas être isolé. De plus, chaque culture ne compte qu’un plot par modalité, sans répétition. 
Ces résultats suggèrent donc que l’amendement a amélioré l’état du sol des plots tests dans les conditions de cet essai, sans qu’ils puissent être généralisés.
